% =====================================================================
\section{ETSI GS QKD 018}
\label{sec:qkd018}
% =====================================================================

La GS~QKD~018 (V1.1.1, aprile 2022)~\cite{qkd018} definisce
l'\emph{orchestration interface} tra un SDN orchestrator e l'SDN controller
di una rete QKD. Il controller agisce da \emph{server}, l'orchestrator da
\emph{client}; il modello informativo è espresso in YANG ed è indipendente
dal vendor. Nella terminologia SDN è la \emph{northbound interface} del
controller QKD, cioè l'esatto complemento della GS~QKD~015. La gerarchia
complessiva diventa:
\begin{center}
\small
orchestrator $\xrightarrow{\text{QKD 018}}$ controller
$\xrightarrow{\text{QKD 015}}$ SDN agent / nodo SD-QKD
$\xrightarrow{\text{QKD 004/014}}$ SAE.
\end{center}

\begin{definition}[Orchestrazione SDN]
L'\emph{orchestrazione SDN} è il processo continuo con cui le risorse
disponibili vengono coordinate automaticamente, secondo criteri di
ottimizzazione, per stabilire e rilasciare servizi end-to-end che
attraversano più domini di rete, ciascuno governato dal proprio controller
SDN~\cite[§4.2]{qkd018}.
\end{definition}

\begin{definition}[Service link]
Un \emph{service link} è un collegamento logico di associazione di chiave
tra due nodi QKD, realizzato da uno o più link fisici QKD (single-hop o
multi-hop). A livello di orchestrazione corrisponde al link virtuale della
GS~QKD~015.
\end{definition}

% Figura: scenario di riferimento della GS QKD 018 (adattato da GS QKD 018, Fig. 1)
\begin{figure}[t]
\centering
\begin{tikzpicture}
  \node[box,fill=blue!8,minimum width=3.3cm] (orch) {SDN orchestrator};
  \node[box,fill=green!8,minimum width=3.3cm,below left=12mm and 6mm of orch] (cq)
       {SDN controller\\rete QKD};
  \node[box,fill=orange!10,minimum width=3.3cm,below right=12mm and 6mm of orch] (co)
       {SDN controller\\rete OTN};
  \node[box,fill=green!4,below=10mm of cq,minimum width=4.6cm] (nq)
       {nodi SD-QKD (agent, KME)};
  \node[box,fill=orange!5,below=10mm of co,minimum width=4.6cm] (no)
       {nodi OTN (SAE)};
  \draw[<->,very thick,red!70!black] (orch) -- node[lbl]{GS QKD 018} (cq);
  \draw[<->,thick] (orch) -- node[lbl]{NBI OTN (fuori scope)} (co);
  \draw[<->,thick] (cq) -- node[lbl]{GS QKD 015} (nq);
  \draw[<->,thick] (co) -- (no);
  \draw[<->,dashed] (nq) -- node[lbl]{QKD 004 / 014} (no);
\end{tikzpicture}
\caption[Scenario di riferimento della GS QKD 018]{Scenario di riferimento della GS~QKD~018 (adattato
da~\cite[Fig.~1]{qkd018}). In rosso l'interfaccia standardizzata; la
linea tratteggiata è la consegna delle chiavi dai KME ai SAE, che avviene
nello stesso perimetro di sicurezza.}
\label{fig:scenario}
\end{figure}


% ---------------------------------------------------------------------
\subsection{Use case della QKD 018}
\label{sec:usecase018}

Il caso d'uso da cui nasce lo standard (§4.1) riprende quasi alla lettera
il problema del backbone della GS~QKD~002. Un operatore che introduce la
QKD può far convivere canali quantistici e classici nella stessa fibra
oppure tenerli separati. La convivenza è difficile per due motivi: le
potenze ottiche dei canali classici sono di molti ordini di grandezza
superiori e disturbano il canale quantistico, e le distanze raggiungibili
dai due tipi di canale sono molto diverse. In un ambiente ricco di fibra
l'operatore può quindi scegliere una rete QKD separata dalla rete di
trasporto ottico (OTN), con vantaggi di prestazioni, isolamento e
separazione delle responsabilità. Seguendo il paradigma SDN, le due reti
sono governate da controller diversi.

Nasce così un problema di coordinamento: i SAE che devono cifrare il
traffico stanno nei nodi OTN, le chiavi sono generate nei nodi QKD e
ciascun controller conosce solo la propria rete. L'orchestrator risolve tre
questioni:

\begin{enumerate}
  \item \textbf{Address matching}: associa l'indirizzo del nodo OTN che deve
        ricevere chiavi a quello del nodo QKD che le fornisce; i due si
        trovano nello stesso perimetro di sicurezza (Fig.~\ref{fig:scenario});
  \item \textbf{provisioning end-to-end}: traduce la richiesta di un cliente
        nella configurazione dei singoli domini, sceglie i domini da
        attraversare e coordina i cammini;
  \item \textbf{gestione}: raccoglie topologia, inventario delle risorse di
        chiave e informazioni FCAPS di ogni dominio.
\end{enumerate}

Generalizzando, lo stesso schema permette di orchestrare \emph{più domini
QKD di vendor diversi}, ciascuno dietro il proprio controller, oltre a
domini QKD e classici insieme. Lo standard sottolinea anche un vantaggio
pratico: un operatore che ha già un orchestrator per le proprie reti
ottiche può integrare la rete QKD senza introdurre strumenti di gestione
separati, che è esattamente il requisito di gestione integrata del caso UC4.

Le funzioni richieste all'orchestrator (§4.2) sono:
\begin{itemize}
  \item tradurre le richieste di servizio dei clienti nella configurazione
        dei domini e nell'allocazione dei SAE;
  \item richiedere attivazione e rilascio dei servizi, incluse le
        connessioni tra domini;
  \item calcolare cammini multi-dominio, con le connessioni inter-dominio e
        gli estremi;
  \item modificare i cammini su vincoli dell'operatore;
  \item scoprire nodi e link di ciascun controller e mantenerne una vista
        astratta;
  \item monitorare le risorse di chiave disponibili su nodi e link;
  \item monitorare la gestione FCAPS e riceverne le notifiche.
\end{itemize}

% ---------------------------------------------------------------------
\subsection{Il modello informativo}
\label{sec:infomodel}

La clausola 5 traduce queste funzioni in cinque blocchi informativi.

\paragraph{Discovery della topologia (§5.1).}
Il controller espone nodi SD-QKD, interfacce e link fisici, e in seguito i
service link. L'orchestrator non ha bisogno della composizione fisica
completa, ma di una vista \emph{astratta} di ciascun dominio e della
connettività tra domini. La discovery può essere \emph{proattiva}
(l'orchestrator interroga il controller prima di ogni nuovo calcolo di
cammino) o \emph{reattiva} (il controller notifica ogni cambiamento, per
esempio l'aggiunta o la rimozione di un nodo).

\paragraph{Stato della rete e inventario (§5.2).}
Per ogni nodo e link fisico il controller espone lo stato, le prestazioni
fisiche e soprattutto le risorse di chiave: il tasso di generazione e il
tasso disponibile dopo il consumo interno. Per i service link si aggiungono
la banda, la chiave disponibile e la politica di ripristino in caso di
guasto.

\paragraph{Stato dei servizi end-to-end (§5.3).}
L'orchestrator può interrogare le applicazioni esterne attive, i nodi
locale e remoto e i service link che le servono, con il relativo stato e
tasso di chiave.

\paragraph{Provisioning con calcolo dei cammini (§5.4).}
L'orchestrator chiede al controller cammini che soddisfino certi vincoli; il
controller restituisce una lista di candidati, eventualmente più volte con
vincoli diversi, finché l'orchestrator non ne sceglie uno da attivare. In
caso di guasto l'orchestrator può chiedere di sostituire il cammino.

\paragraph{Notifiche (§5.5).}
Il controller segnala eventi, allarmi e superamenti di soglia. Le soglie
sulle prestazioni (QBER, SKR, ESKR) sono configurabili dall'orchestrator.
Rispetto alla 015 si aggiungono notifiche sullo stato dei nodi
(\texttt{sdqkdn\_new}, \texttt{\_operating}, \texttt{\_down},
\texttt{\_failure}, \texttt{\_out}), sulle applicazioni attivate
(\texttt{sdqkdn\_application\_deployed}) e sui service link
(\texttt{sdqkdn\_svc\_link\_down}, \texttt{\_perf\_update},
\texttt{\_overloaded}).

La~\eqref{eq:eskr015} si estende in modo naturale al livello di
orchestrazione. Poiché in un cammino multi-hop $p$ ogni link inoltra la
stessa chiave, un nuovo servizio con richiesta $r$ è sostenibile solo se il
collo di bottiglia lo regge:
\begin{equation}
  p \text{ ammissibile} \iff
  \min_{\ell\in p}\, \mathrm{ESKR}_\ell \;\ge\; r .
  \label{eq:admissible}
\end{equation}
Anche questa è una formalizzazione nostra, coerente con i campi
\texttt{skr}, \texttt{eskr} ed \texttt{expected\_consumption} dei modelli
YANG; lo standard non prescrive alcun algoritmo di calcolo dei cammini.
